
\section{AT-05: Hartree--Fock direct and exchange matrix elements}
\source{2026 Atomic Assignment 4 (\texttt{PHYS4055\_A2-3.pdf}), Question 1.}
\subsection*{Step 1: antisymmetrise two occupied spin-orbitals}
Let \(a(x)\), \(b(x)\) be orthonormal spin-orbitals;
\(x=(\mathbf r,\sigma)\). A two-electron determinant is
\[
\Psi_{ab}(1,2)=\frac1{\sqrt2}
\left[a(1)b(2)-b(1)a(2)\right].
\]
The electron repulsion is \(v_{12}=e^2/(4\pi\epsilon_0r_{12})\).
Expanding \(\langle\Psi_{ab}|v_{12}|\Psi_{ab}\rangle\)
gives four terms. Interchanging labels \(1\leftrightarrow2\)
shows that the two direct terms agree and the two exchange
terms agree:
\[
\boxed{\langle v_{12}\rangle=J_{ab}-K_{ab}}.
\]
With atomic units, these integrals are
\[
\begin{aligned}
J_{ab}&=\iint
\frac{|a(x_1)|^2|b(x_2)|^2}{r_{12}}\,dx_1dx_2,\\
K_{ab}&=\iint
\frac{a^*(x_1)b(x_1)b^*(x_2)a(x_2)}{r_{12}}
\,dx_1dx_2.
\end{aligned}
\]
The exchange term comes from fermionic antisymmetry;
it is not a separate classical Coulomb force.

\subsection*{Step 2: integrate spins before orbital angles}
For separable \(a(x)=\phi_a(\mathbf r)\chi_a(\sigma)\),
\(b(x)=\phi_b(\mathbf r)\chi_b(\sigma)\), the direct term
is independent of the relative spin. The exchange term
contains \(|\langle\chi_a|\chi_b\rangle|^2\).
For orthogonal spin states \(\alpha\) and \(\beta\),
this factor is zero: \(\boxed{K_{a\alpha,b\beta}=0}\).
For equal spin states, \(K_{ab}\) generally survives.

\subsection*{Step 3: reduce an illustrative \(s\)-orbital pair}
Take orbitals with \(l_a=l_b=0\) and
\(\phi_a(\mathbf r)=P_a(r)Y_{00}(\Omega)/r\)
with \(\int_0^\infty |P_a(r)|^2dr=1\).
In the multipole expansion of \(1/r_{12}\),
only \(k=0\) survives the angular integrations.
Write \(r_>=\max(r_1,r_2)\). Then
\[
\boxed{J_{ab}=\int_0^\infty\!\int_0^\infty
\frac{|P_a(r_1)|^2|P_b(r_2)|^2}{r_>}
\,dr_1dr_2}
\]
and for equal spin
\[
\boxed{K_{ab}=\int_0^\infty\!\int_0^\infty
\frac{P_a^*(r_1)P_b(r_1)P_b^*(r_2)P_a(r_2)}
{r_>}\,dr_1dr_2}.
\]
For higher \(l\), retain the appropriate
\(k,q\) angular coefficients (Gaunt/3-j factors).

\subsection*{Interpretation}
The direct mean field is the Coulomb potential
of the other electron density; exchange is a
nonlocal, spin-selective operator. In Hartree--Fock,
optimising the determinant self-consistently produces
both terms in the one-electron equations.
\checkmethod{The full Assignment 4 derivation requires
the Johnson-textbook angular and radial conventions.
This example isolates the physical and algebraic
direct--exchange mechanism, not every general \(k\) term.}
